% Day 3 frames moved from Day 1 on 2026-10-17 (Day 1 had 44 frames). \input this from day3.tex.
% Figures d1_candidates.pdf / d1_diminishing.pdf are still produced by figures/make_day1.py.

% =================================================================================================
\section{Why action needs a second functional}
% =================================================================================================

\begin{frame}{Perception is easy to state. Action is not.}
  \textbf{Perception.} $o_t$ is fixed, so only $q$ can move:
  \[ q^\ast = \argmin_q F[q,o_t] \]
  That is approximate Bayesian inference, and nothing else.

  \pause
  \vspace{0.6em}
  \textbf{Action.} At decision time $o_t$ is a number in a buffer. No action can change it:
  \[ \frac{\partial F[q,o_t]}{\partial a} = 0 \]
  An agent descending the instantaneous free energy finds a \alert{flat landscape} in every
  direction.

  \pause
  \vspace{0.6em}
  What action \emph{can} change is the free energy of observations not yet received:
  \[ a^\ast = \argmin_a \E_{q(o_{t+1},s_{t+1}\given a)}\big[F[q,o_{t+1}]\big] \]
  This is not a new principle. It changes what is held fixed and what is averaged over.
  \src{active\_inference.tex l.335--369}
\end{frame}

\begin{frame}{Expected free energy}
  For a policy $\pi$ (a sequence of actions), with the predicted states and observations
  $\tilde s,\tilde o$:
  \[ G(\pi) = \E_{q(\tilde s,\tilde o\given\pi)}\big[\ln q(\tilde s\given\pi) - \ln p(\tilde o,\tilde s)\big] \]
  Goals enter as a \textbf{prior over preferred outcomes} $p(\tilde o)$. There is no separate
  reward channel. Then, exactly:
  \[
    G(\pi) = \E\Big[
      \underbrace{\ln q(\tilde s\given\pi) - \ln q(\tilde s\given\tilde o,\pi)}_{\text{epistemic}}
      \;\underbrace{-\ln p(\tilde o)}_{\text{pragmatic}}
      + \underbrace{\ln q(\tilde s\given\tilde o,\pi)-\ln p(\tilde s\given\tilde o)}_{\to 0 \text{ if inference stays good}}
    \Big]
  \]
  \src{active\_inference.tex l.371--412}
\end{frame}

\begin{frame}{Two terms, one currency}
  \begin{columns}[T]
    \begin{column}{0.48\textwidth}
      \begin{block}{Epistemic value}
        $=-I(\tilde s;\tilde o\given\pi)$: minus the mutual information between future states and
        future observations.\\[0.3em]
        Largest when what you would see depends strongly on what is true. Curiosity is a
        \emph{term of the objective}, not a heuristic bolted on.
      \end{block}
    \end{column}
    \begin{column}{0.48\textwidth}
      \begin{block}{Pragmatic value}
        $=-\E\ln p(\tilde o)$: how improbable the predicted outcomes are under the agent's
        preferences.\\[0.3em]
        This is goal-seeking: the part a classical cost function would have captured.
      \end{block}
    \end{column}
  \end{columns}
  \vspace{0.6em}
  Both are in \textbf{nats}, and they were derived together. The exploration/exploitation balance is
  set by the algebra, not by a mixing coefficient. It moves by itself as the beliefs sharpen.
  \src{active\_inference.tex l.414--427}
\end{frame}

\begin{frame}{For Gaussian beliefs, three names for one number}
  Prior precision $\Lambda$, a measurement with Jacobian $J$ and noise $R$:
  \[
    \underbrace{H[q]-\E H[q\given o]}_{\text{expected entropy drop}}
    = \underbrace{I(s;o)}_{\text{mutual information}}
    = \underbrace{\tfrac12\ln\frac{|\Lambda + J^\top R^{-1}J|}{|\Lambda|}}_{\text{D-optimal design}}
  \]
  \begin{itemize}
    \item The epistemic term of active inference, the information gain of active sensing, and
          the D-criterion of experiment design are \alert{the same quantity}.
    \item It depends on \emph{where} you look ($J$), not on \emph{what} you will see. So it can be
          computed before moving.
  \end{itemize}
  \src{appendix\_math.tex l.742--779, l.1236--1287}
\end{frame}

\begin{frame}{Where to look next}
  \begin{columns}[c]
    \begin{column}{0.58\textwidth}
      \includegraphics[width=\textwidth]{d1_candidates.pdf}
    \end{column}
    \begin{column}{0.4\textwidth}
      \small
      The robot (red) has seen the east part of the north wall. Five candidate viewpoints:
      \begin{itemize}
        \item ``far west'' sees the \emph{fewest} new points except one, yet gains the
              \emph{most}: it adds lever arm where the current information is weak;
        \item ``stay'' gains exactly $\ln 2$. Why?
      \end{itemize}
    \end{column}
  \end{columns}
\end{frame}

\begin{frame}{Diminishing returns, and a warning}
  \begin{columns}[c]
    \begin{column}{0.45\textwidth}
      \includegraphics[width=\textwidth]{d1_diminishing.pdf}
    \end{column}
    \begin{column}{0.53\textwidth}
      \small
      The $k$-th identical look gives $|\,(k{+}1)\Lambda_1| / |\,k\Lambda_1|=((k{+}1)/k)^2$ for
      two parameters, so
      \[ \Delta H_k = \ln\frac{k+1}{k}. \]
      \pause
      \alert{But this assumes each look's errors are independent.} On a robot, repeated looks
      share the pose error, the mount error and the calibration error, and those do not average
      away.\\[0.4em]
      Day 2's common-mode ceiling already put a number on how little a repeat look is worth.
    \end{column}
  \end{columns}
\end{frame}

\begin{frame}{Pricing travel in the same currency}
  The simplest two-term expected free energy, the one that drives the robot's viewpoint selection:
  \[
    G(v) = \underbrace{\beta\,\lVert v - x\rVert}_{\text{cost of travel}}
         - \underbrace{\Delta H(v)}_{\text{information gain}}
  \]
  \begin{itemize}
    \item $\beta$ is an \emph{exchange rate}, in nats per metre, not a tuning knob between two
          incommensurable objectives.
    \item In the lab: above $\beta\approx0.07$ nats/m the robot stays put and scans again,
          because the (overvalued) $\ln 2$ is free.
    \item The full version sums agent-owned terms over a horizon. That is Day 5.
  \end{itemize}
  \src{introduction.tex l.244--253; action.tex}
\end{frame}

\begin{frame}{One objective across time scales}
  The gradient of a $T$-step expected free energy with respect to the current command splits as
  \[
    \nabla_{u_t} G_T(\pi)
      = \underbrace{\nabla_{u_t} G_{t+1}(u_t)}_{\text{what a reactive controller computes}}
      + \underbrace{\sum_{\tau=t+2}^{t+T}\nabla_{u_t} G_\tau(\pi)}_{\text{what planning adds}}
  \]
  \begin{itemize}
    \item Reactive control is the leading term of planning. They are not two systems.
    \item The relatives you know (planning-as-inference, MPC, MPPI) sit on this line. They differ in
          how many terms they keep and how they estimate them.
  \end{itemize}
  \src{active\_inference.tex l.431--510, l.557--579}
\end{frame}

% --- measured results that need expected free energy ------------------------------------------

\begin{frame}{Is the derived objective better than a good heuristic?}
  The room explorer, written in nats with no free weights, against the hand-tuned scorer it
  replaced. Both were run over the \textbf{same 594 real room layouts}, paired.
  \vspace{0.3em}
  \begin{center}\small
  \begin{tabular}{@{}l r r r@{}}
    \toprule
     & IoU mean & IoU median & Hausdorff median \\
    \midrule
    hand-tuned scorer & 0.9794 & 0.9880 & 0.038 m \\
    derived objective & 0.9806 & 0.9870 & 0.069 m \\
    \midrule
    paired difference & \multicolumn{3}{l}{$+0.0013\pm0.0030$; better / worse / tied $=238/298/58$, sign test $p=0.011$} \\
    \bottomrule
  \end{tabular}
  \end{center}
  \begin{itemize}\small
    \item It wins the mean and loses the count. Five rescued rooms carry the whole mean; without
          them the difference is $-0.00085$.
    \item The verdict is \alert{equivalence}, and the case for the derived form rests on its having
          no free parameters, not on performance.
  \end{itemize}
  \small On Day 5 we will see that even this IoU was not measuring what we thought it was.
  \src{experiments.tex l.1404--1446}
\end{frame}

\begin{frame}{Doing nothing can be an epistemic act}
  \begin{measured}{gyro bias, two legs of the same tour}
    The share of the bias posterior that came from data rather than from the prior: about
    \textbf{half} on one leg, \textbf{almost all} on the other. The only difference was that the
    robot happened to stand still for a few minutes.
  \end{measured}
  \begin{itemize}
    \item The bias regresses on elapsed time alone, so the action that teaches it is the one that
          only spends time: \alert{stopping}.
    \item The body is instrumented by everything the robot does, because every action passes
          through it. The world is sampled only where the robot looks.
    \item Hence: epistemic actions for rooms and objects, almost none for the body (Part A).
  \end{itemize}
  \src{active\_inference.tex l.850--880, l.952--968}
\end{frame}

